% Architecture des modules : graphe de dépendance d'inclusion.
\begin{tikzpicture}[
  node distance = 1.4cm and 1.6cm,
  module/.style = {
    draw, rounded corners, thick,
    minimum width = 2.2cm, minimum height = 0.9cm,
    align = center, font = \small\ttfamily,
    fill = blue!8
  },
  vendor/.style = {
    draw, rounded corners, thick, dashed,
    minimum width = 2.2cm, minimum height = 0.9cm,
    align = center, font = \small\ttfamily,
    fill = orange!10
  },
  arr/.style = {-Stealth, thick, gray!70}
]

% Niveau 1 : main
\node[module] (main) {main};

% Niveau 2 : world, conf
\node[module, below left = of main]  (world) {world};
\node[module, below right = of main] (conf)  {conf};

% Niveau 3 : entity, species
\node[module, below = of world] (entity)  {entity};
\node[module, below = of conf]  (species) {species};

% Niveau 4 : vec2, tomlc99
\node[module, below = of entity] (vec2) {vec2};
\node[vendor, below = of species] (toml) {tomlc99\\(vendor)};

% Inclusions
\draw[arr] (main)    -- (world);
\draw[arr] (main)    -- (conf);
\draw[arr] (world)   -- (entity);
\draw[arr] (world)   -- (species);
\draw[arr] (entity)  -- (species);
\draw[arr] (entity)  -- (vec2);
\draw[arr] (world)   -- (vec2);
\draw[arr] (conf)    -- (species);
\draw[arr] (conf)    -- (toml);

% Légende
\node[draw=none, font=\footnotesize, below = 1cm of vec2,
      anchor=north west, xshift=-0.5cm] (legend) {
  \begin{tabular}{l}
    \tikz \draw[-Stealth, thick, gray!70] (0,0) -- (0.5,0); \ \emph{inclusion}\\[2pt]
  \end{tabular}
};

\end{tikzpicture}
